\documentclass[12pt]{article}
\usepackage{graphicx} % Required for inserting images

\title{NPRE 455 Computer Project 1}
\author{Owen Strong \\
\itshape Prof.: April Novak\\
% \itshape TA: Kholod Mahmoud \\
% \itshape TA: Shaffer Bauer \\
% \itshape TA: Justin Jia \\
}
\date{March 24th, 2026}

\usepackage[
    style=ieee
]{biblatex}
\bibliography{refs.bib}

\usepackage[letterpaper, margin=1in]{geometry}
\usepackage{xspace}

% Figure Macros
\newcommand{\figwidth}{0.75\linewidth}
\newcommand{\figref}[1]{Fig.~\ref{#1}}
\newcommand{\defplotW}[3]{
    \begin{figure}[!htb]
        \centering
        \includegraphics[width=#3\linewidth]{fig/#1.png}
        \caption{#2}
        \newcommand{\tmplabel}{#1}
        \label{\tmplabel}
    \end{figure}
}
\newcommand{\defplot}[2]{\defplotW{#1}{#2}{0.75}}

% Common Macros
\newcommand{\micCi}{\xspace$\mathrm{\mu Ci}$\xspace}
\newcommand{\CsSrc}{\xspace$\mathrm{^{137}Cs}$\xspace}
\newcommand{\bigO}{\mathcal{O}}

\usepackage{placeins} % FloatBarrier
\usepackage{amsmath} % align
\usepackage{xcolor}
\usepackage[colorlinks=true, citecolor=violet, linkcolor=violet, urlcolor=violet]{hyperref}
\usepackage{enumitem}

\begin{document}
\pagenumbering{gobble}
\maketitle
\pagebreak
{\hypersetup{linkcolor=black}\tableofcontents}
\pagebreak
\pagenumbering{arabic}
\begin{abstract}
    The Neutron Diffusion Equation is a key model for understanding the behavior of neutrons in a system like a nuclear reactor, and is often approximated using methods such as finite-difference approximation to model more complicated geometry than would be feasible using analytic solutions alone.
    In this report, we aim to explore finite-difference methods and their application to the Neutron Diffusion Equation, namely to multiplying and non-multiplying planar systems.
    We find finite difference approximations of a non-multiplying and a multiplying form of the Diffusion Equation, and apply them to models of the planar systems discretized in space over various sizes $h$ of cells.
    We find a reasonable numerical solution convergence of order $\bigO(h^{1.009})$ in the non-multiplying system to the analytical solution and a consistent convergence for different combinations of geometry for the multiplying system to the flux distributions and multiplication values found using analytical solutions.
    These results confirm the effectiveness and certain key limitations of finite-difference approximations of the Neutron Diffusion Equation, and may be of value to those aiming to perform neutronics analysis with limited time and resources.
\end{abstract}

\section{Introduction}
Neutronics, or the behavior of neutrons in a machine such as a nuclear reactor, is a crucial field of study in nuclear engineering.
Nuclear reactor design, operation, decommissioning, and safety every step of the way all depend on being able to predict and understand the movement and interaction of neutrons with matter inside and out of the reactor.\\

While several approaches are available towards characterizing neutron behavior in materials, perhaps the most common starting point is the robust Neutron Transport Equation, and the simpler, derived-therefrom, Neutron Diffusion Equation. 
These equations are either solved analytically (when feasible, typically only for the Diffusion Equation in limited geometry), or more conventionally, using numerical methods like finite-difference approximations, to handle more complicated geometry and other requirements.\\

In this report, finite-difference approximations are explored briefly in general, and their application to solve two one-dimensional diffusion problems (one of a distributed source region, and one of a multiplying region) tested and compared to analytical solutions to find flux distributions, and for the multiplying system, multiplication rates $k$.
\section{Problem 1: Finite Difference Overview}
% Theory behind finite difference methods
In engineering mathematics, a number of differential equations prove undesirable to solve using analytic techniques.
A common practical technique to solve problems using these equations is instead to use a ``finite difference'' approximation scheme, in which derivatives are approximated using finite differential quotients.\\

\subsection{Finite-Difference Theory}
Consider, in the simplest case, the definition of a 1D derivative as the slope, or rise over run, of a function:
\begin{equation}
    \label{eq:lim-h}
    f'(x)=\frac{df(x)}{dx}=\lim_{h\rightarrow0}\frac{f(x+h)-f(x)}{h}
\end{equation}
This yields a simple ``finite difference'' approximation of the derivative, with a small enough spacing $h$:
\begin{equation}\label{eq:right}
    f'(x)\approx\frac{f(x+h)-f(x)}{h}
\end{equation}
There is, however, no reason this could not be taken about a point $h$ behind $x$:
\begin{equation}
    f'(x)\approx\frac{f(x)-f(x-h)}{h}
\end{equation}
or from both sides:
\begin{equation}\label{eq:central}
    f'(x)\approx\frac{f(x+h)-f(x-h)}{2h}
\end{equation}
To question how accurate such an approximation is, one can employ a Taylor series approximation of the function around $x$:
\begin{equation}
    f(x+h)=f(x)+hf'(x)+\frac{h^2}{2}f''(x)+\frac{h^3}{6}f'''(x)+HOT
\end{equation}
where $HOT$ denotes ``higher order terms'' which we expect to contribute increasingly less and less.
In this case, by applying the approximation at $x+h$, we can derive the right side differential from Eq.~\ref{eq:right}:
\begin{gather*}
    f(x+h)=f(x)+hf'(x)+\frac{h^2}{2}f''(x)+\frac{h^3}{6}f'''(x)+HOT \\
    f(x+h)-f(x)=hf'(x)+\frac{h^2}{2}f''(x)+\frac{h^3}{6}f'''(x)+HOT
\end{gather*}
\begin{equation}\label{eq:right_err}
    \frac{f(x+h)-f(x)}{h}=f'(x)+\boxed{\frac{h}{2}f''(x)+\frac{h^2}{6}f'''(x)+HOT}
\end{equation}
where the boxed term now quantifies the expected difference of Eq~\ref{eq:right}'s approximation from the true $f'(x)$.
If a small enough $h$ is chosen ($h<1$), we expect $h^2$ to decay faster than $h$ as $h$ decreases, and $h^3$ than $h^2$ and so on.
Thus, since the lowest power of $h$ an error term scales with in Eq.~\ref{eq:right_err} is $h$, we say that this approximation has an error on the order $\bigO(h)$.\\

This use of the Taylor series, of course, can be applied to other points than $x+h$:
\begin{gather*}
    f(x-h)=f(x)-hf'(x)+\frac{h^2}{2}f''(x)-\frac{h^3}{6}f'''(x)+HOT \\
    f(x)-f(x-h)=hf'(x)-\frac{h^2}{2}f''(x)+\frac{h^3}{6}f'''(x)+HOT \\
\end{gather*}
\begin{equation}
    \frac{f(x)-f(x-h)}{h}=f'(x)-\frac{h}{2}f''(x)+\frac{h^2}{6}f'''(x)+HOT
\end{equation}
or combined as in Eq.~\ref{eq:central}:
\begin{align*}
    f(x+h)-f(x-h)=&f(x)+hf'(x)+\frac{h^2}{2}f''(x)+\frac{h^3}{6}f'''(x)+HOT\\
    &-\left[f(x)-hf'(x)+\frac{h^2}{2}f''(x)-\frac{h^3}{6}f'''(x)+HOT\right]\\
    =&2hf'(x)+2\frac{h^3}{6}f'''(x)+HOT
\end{align*}
\begin{equation}\label{eq:central_err}
    \frac{f(x+h)-f(x-h)}{2h}=f'(x)+\frac{h^2}{6}f'''(x)+HOT
\end{equation}
Note how in this last example, a pair of terms cancel out, resulting in an error on the order $\bigO(h^2)$ instead of $\bigO(h)$.\\

In fact, given some number $n$ of accessible \emph{nodes} of $f(x)$, one can find \emph{weights} for a generally optimal linear combination of these points equivalent to $f'(x)$ which eliminates at least $n-1$ error terms.

\subsection{Right-handed Finite-Difference}
For example, consider using the points $f(x)$, $f(x+h)$, $f(x+2h)$, and $f(x+3h)$. %, denoting these $x$ values $x_i$ through $x_{i+3}$.
If we rewrite these points as Taylor expansions around $x$:
\begin{align}
    f(x) &= f(x) && + 0 && + 0 && + 0 && + ...\\
    f(x+h) &= f(x) &&+ hf'(x) &&+ \frac{h^2}{2}f''(x) &&+ \frac{h^3}{6}f'''(x) &&+...\\
    f(x+2h) &= f(x) &&+ 2hf'(x) &&+ \frac{4h^2}{2}f''(x) &&+ \frac{8h^3}{6}f'''(x) &&+...\\
    f(x+3h) &= f(x) &&+ 3hf'(x) &&+ \frac{9h^2}{2}f''(x) &&+ \frac{27h^3}{6}f'''(x) &&+...\\
\end{align}
thus, having four expressions to combine there is some linear combination of these points $f(x+nh)$ that both is equivalent to the target derivative and eliminates the first three other terms of the polynomials, leaving only error from the higher order terms of maximum order $\bigO(h^4)$.
In fact, we could choose any of these first three derivatives to feasibly approximate.
Suppose we wished to find some linear combination $af(x) + bf(x+h)+cf(x+2h)+df(x+3f)$ which approximated $f''(x)$ with as small an order of error as possible.\footnote{Since the referenced points are all to the right of $x$, this would be considered a \emph{right}-handed, or \emph{forward} finite difference approximation. One entirely on the left side would be considered left-handed or backwards, while a perfectly centered one such as $af(x-h)+bf(x)+cf(x+h)$ would be considered a central difference.}
Expanding this combination (and introducing the notation $x_{i+n}=x+nh$):
\begin{align}
    af(x_i) + bf(x_{i+1})+cf(x_{i+2})+df(x_{i+3})
    &=(a+b+c+d)f(x_i)\\
    &+(b + 2c + 3d)hf'(x_i)\nonumber\\
    &+(b + 4c + 9d)h^2f''(x_i)/2\nonumber\\
    &+(b + 8c + 27d)h^3f'''(x_i)/6\nonumber\\
    &+(b + 16c + 81d)h^4f'''(x_i)/24+...\nonumber
\end{align}
We can formulate this as a matrix problem to have $1\times f''(x_i)$, and eliminate the first three other terms.
That is:
\begin{equation}\label{eq:find_coeff}
    \begin{bmatrix}
        1 & 1 & 1 & 1 \\
        0 & h & 2h & 3h \\
        0 & h^2/2 & 4h^2/2 & 9h^2/2 \\
        0 & h^3/6 & 8h^3/6 & 27h^3/6 \\
    \end{bmatrix}
    \begin{bmatrix}
        a\\b\\c\\d
    \end{bmatrix}
    =
    \begin{bmatrix}
        0\\0\\1\\0
    \end{bmatrix}
\end{equation}
to achieve:
\begin{align}
    af(x_i) + bf(x_{i+1})+cf(x_{i+2})+df(x_{i+3})
    &=(0)f(x_i)\\
    &+(0)f'(x_i)\nonumber\\
    &+(1)f''(x_i)\nonumber\\
    &+(0)f'''(x_i)\nonumber\\
    &+(b + 16c + 81d)h^4f''''(x_i)/24+...\nonumber
\end{align}
Solving the problem of Eq.~\ref{eq:find_coeff} yields:
\begin{equation}
    \begin{bmatrix}
        a\\b\\c\\d
    \end{bmatrix}
    =
    \frac{1}{h^2}
    \begin{bmatrix}
        1\\-5\\4\\-1
    \end{bmatrix}
\end{equation}
yielding a final approximation:
\begin{align}\label{eq:final_werror}
    \frac{1f(x_i) - 5f(x_{i+1})+4f(x_{i+2})-f(x_{i+3})}{h^2}&=f''(x_i)+(\frac{-5+4\times16-1\times81}{24})h^2f''''(x_i)+HOT\nonumber\\
    \frac{1f(x_i) - 5f(x_{i+1})+4f(x_{i+2})-f(x_{i+3})}{h^2}&=f''(x_i)+\frac{-22}{24}h^2f''''(x_i)+HOT
\end{align}
or, for a function $\phi(x)$, to estimate $\frac{d\phi(x_i)}{dx}$ with $\phi(x_i)$, $\phi(x_{i+1})$, $\phi(x_{i+2})$, $\phi(x_{i+3})$:
\begin{equation}
    \frac{d\phi(x_i)}{dx}\approx\frac{1\phi(x_i) - 5\phi(x_{i+1})+4\phi(x_{i+2})-1\phi(x_{i+3})}{h^2}
\end{equation}
As seen in Eq.~\ref{eq:final_werror}, the most substantial term we were unable to eliminate was $\frac{-22}{24}h^2f''''(x)$, giving this approximation an error on the order $\bigO(h^2)$.\\

To put this another way, suppose that one wished to refine $h$ to achieve an error reduction of $10\times$.
If an initial approximation with spacing $h_0$ yielded error on the order $\bigO(h_0^2)$, then one could choose $h_1=Ch_0$:
\begin{align}
    \bigO(h_0^2)/10&=\bigO(h_1^2)\\
    &=\bigO((Ch_0)^2)\nonumber\\
    \bigO(h_0^2)/10&=C^2\bigO(h_0^2)\nonumber\\
    C&=\sqrt{1/10}
\end{align}
yielding $h_1=h_0\sqrt{\frac{1}{10}}$.\\

\figref{phi2p} illustrates this approximation used for $\phi(x)=sin(x)$ using different values of $h\in\{1, 0.5, 0.25\}$, and compares to an analytical solution.
Note how while the solution using the coarsest $h=1$ is extremely rough, every halving of $h$ drastically improves the accuracy of the approximation, with $h=0.25$ already resulting in a very believable approximation.
\defplot{phi2p}{The second derivative of $\phi(x)=sin(x)$, as found using a finite difference approximation over differing fineness of $h$, compared to the analytical solution $\phi''(x)=-sin(x)$}

\section{Problem 2: Non-multiplying Systems}
\FloatBarrier
\newcommand{\angFlux}{\psi(\vec{r}, E, \hat{\Omega}, t)}
\newcommand{\sclFlux}{\phi(\vec{r}, E, t)}
The aggregate movement of neutrons can be described using the \emph{Neutron Transport Equation}, a balance equation describing the population of neutrons (as angular $\angFlux$ and scalar $\sclFlux$ neutron flux) at a given point, energy, direction, and time:
\begin{align}
    \frac{1}{v(E)}\frac{\partial\angFlux}{\partial t} &\qquad \mathrm{Change~in~time}\label{eq:trans_dt}\\ 
    +\hat{\Omega}\cdot\nabla\angFlux &\qquad \mathrm{Neutron~streaming~loss}\label{eq:trans_stream}\\
    + \Sigma_t(\vec{r}, E, t)\angFlux &\qquad \mathrm{Collision~loss}\label{eq:trans_coll}\\
    =&\nonumber\\
    \int_{4\pi}d\hat{\Omega}\int_0^\infty dE'\Sigma_s(\vec{r}, E'\rightarrow E, \hat{\Omega}'\rightarrow\hat{\Omega},t)\angFlux &\qquad \mathrm{Scattering~to~energy}~E\label{eq:trans_scat}\\
    +\frac{\chi_p(E)}{4\pi}\int_0^\infty dE'\nu_p(E')\Sigma_f(\vec{r},E',t)\sclFlux&\qquad \mathrm{Prompt~fission~neutrons}\label{trans_prompt}\\
    +\sum_{i=1}^{N}\frac{\chi_{d,i}(E)}{4\pi}\lambda_iC_i(\vec{r},t)&\qquad \mathrm{Delayed~fission~neutrons,~}N~\mathrm{groups}\label{eq:trans_delay}\\
    +s(\vec{r},E,\hat{Omega},t)&\qquad\mathrm{External~sources}\label{eq:trans_ext}
\end{align}
However, the angular streaming term introduces complications in actually trying to solve this application.
A more typical representation of this balance is the \emph{Neutron Diffusion Equation}.
By assuming that neutron transport follows a \emph{diffusive} process, so that neutron current $\vec{J}=-D\nabla\phi$, streaming can be written more directly as a function of the scalar flux $\sclFlux$ and the balance can be written simply with respect to $\phi$:
% \newcommand{\phirt}{\phi(\vec{r}, t)}
\newcommand{\rt}{(\vec{r}, t)}
\begin{equation}
    \label{eq:diffusion}
    \frac{1}{v}\frac{\partial \phi\rt}{\partial t}-\nabla\cdot D\nabla\phi\rt + \Sigma_a\rt\phi\rt = \nu_p\Sigma_f\rt\phi\rt + \sum_{i=1}^{N}\lambda_iC_i\rt + s\rt
\end{equation}
The formulation of the diffusion equation in Eq.~\ref{eq:diffusion} also assumes neutrons are all the same speed and not split into multiple ``groups'' of energy $E$ (the one-group model).
This assumption also serves to simplify calculations, but is not as crucial as the diffusion assumption, and the equation could still be split back up into multiple energy groups for a more accurate representation.\\

In a more precise derivation, the neutron diffusion equation can be found by taking two moments, an integral over angle (zeroth moment) and an integral of that function times angle, over angle (first moment), then introducing the assumption that flux is \emph{linearly anisotropic}. This is by far the greatest assumption which simplifies the neutron transport equation to the neutron diffusion equation: that is, while not necessarily perfectly uniform over angle, the flux can at least be described by a linear function of angle.
This report focuses primarily on the application of the end diffusion equation, but more thorough derivations are available.\cite{novak_lecture6_2026}
For the purposes of this report, the following limitations are most crucial:
\begin{itemize}
    \item The diffusion equation fails if there is no medium to scatter in, like a vacuum.
    \item Near strong sources or sinks of neutrons, there is a strong preferential direction for neutrons which breaks the required linear anisotropy.
    \item Strong spatial variations in general tend to introduce stronger angular dependence than our assumption permits.
    \item Caution should be exercised near (within several mean free paths of) problematic boundaries or discontinuities, such as a vacuum, where neutrons that scatter in that direction wouldn't come back.
\end{itemize}

Even with these limitations, the neutron diffusion equation can still be a very effective model of scalar neutron flux.
A typical process to solve neutronics problems analytically using this equation is to split the problem into regions of uniform or simply-distributed properties.
Over each region $i$, a flux $\phi_i$ is then expected (if reasonable) to follow the neutron diffusion equation.
For a steady state problem, in which the time derivative (\ref{eq:trans_dt}) can be neglected, there remains only the second-degree spatial derivative of $\phi$, requiring two known boundary conditions to analytically resolve that solution.
Between regions $a$ and $b$ which meet at $S_{ab}$, we can apply conservation based boundary conditions:
\begin{itemize}
    \item Flux described by both regions' $\phi_i\rt$ should agree at their interface: 
    \begin{equation}
        \phi_a\rt=\phi_b\rt,{\hat{r}\in S_{ab}}
    \end{equation}
    \item The current crossing the boundary $\vec{J}\cdot\hat{n}$ is conserved, accounting for any known neutron production localized entirely within the boundary: 
    \begin{equation}
        \lim_{\vec{r}\rightarrow\vec{r_s}}\left[\hat{n}\cdot\left(\hat{J_2}\rt-\hat{J_1}\rt\right)\right]
    \end{equation}
\end{itemize}
At vacuum boundaries which don't face other neutron-diffusing materials, one could assume incoming current at the boundary is zero:
\begin{equation}
    J_-(\vec{r_s})=\frac{\phi\vec{r_s}}{4}+\frac{D}{2}\nabla\phi(\vec{r_s})\cdot\hat{n}=0
\end{equation}
Which is in practice often shifted to an extrapolated boundary condition \cite{novak_lecture9_2026}:
\begin{equation}
    \phi(\vec{r_s}+2D\hat{n})=0
\end{equation}
Finally, one more common boundary condition is symmetry around some point, line, or plane.
The gradient of flux can be assigned corresponding to a localized source embedded in this feature of symmetry, or without such a local source, simply assumed to be zero:
\begin{equation}
    \nabla\phi(\vec{r_s})=0
\end{equation}
\\

Using some combination of these boundary conditions (two per flux region), an analytic solution for flux over these regions may be obtained.\\




For a non-multiplying system (with no fission), $\Sigma_t$ in \ref{eq:trans_coll} simplifies to simply absorption $\Sigma_a$, and the prompt and delayed fission sources \ref{trans_prompt} and \ref{eq:trans_delay} can be ignored.
For a system in steady state, we can also ignore the time dependence in \ref{eq:diffusion}:
\begin{equation}
    \label{eq:nonmult-diff}
    -\nabla\cdot D\nabla\phi\rt + \Sigma_a\rt\phi\rt = s\rt
\end{equation}

\newcommand{\ddx}{\frac{\partial}{\partial x}}
\newcommand{\ddxx}{\frac{\partial^2}{\partial x^2}}
\FloatBarrier

\defplot{prob2_setup}{The (symmetric) geometry of Problem~2: A central source, then one slab of Material~A, then a slab of Material~B before vacuum. Taken from Computer Project 1 Manual.\cite{novak_computer_2026}}
Consider the 1D system in \figref{prob2_setup}, with infinite slabs of two materials mirrored around a central plane $x=0$. For this system, one may consider it to have dimensions $a=10$~[cm], $b=20$~[cm], region~1 to have material properties $D_1=1$~[cm], $L_1=10$~[cm], and region~1 material properties $D_2=1.7$~[cm], $L_2=5$~[cm].

\subsection{Analytical Solution}
To analytically solve for this distribution over $0\leq x\leq a+b$, one may first split the problem into two equations, one for each materially uniform region.
For region~1, with uniform properties, with a uniform distributed source S:
\begin{align}
    -\ddx D_1\ddx\phi_1 + \Sigma_{a1}\phi_1 &= S \nonumber\\
    \ddxx\phi_1 - \frac{\Sigma_{a1}}{D_1}\phi_1 &= -\frac{S}{D_1} \nonumber\\
    \label{eq:reg1-source}
    \ddxx\phi_1 - \frac{1}{L_1^2}\phi_1 &= -\frac{S}{D_1}
\end{align}
and got region 2, without such a source:
\begin{equation}
    \label{eq:reg2}
    \ddxx\phi_2 - \frac{1}{L_2^2}\phi_2 = 0
\end{equation}
First, consider the region~1 flux $\phi_1$.
Eq.~\ref{eq:reg1-source} is not homogeneous (having the term $\frac{S}{D_1}$ independent of flux $\phi_1$), but can be split into a homogeneous component solution $\phi_h$ and a particular solution $\phi_p$
\begin{equation}
    \label{eq:homo-part}
    \phi_1 = \phi_h + \phi_p
\end{equation}
\begin{align}
    \label{eq:homo}
    \ddxx \phi_h - \frac{1}{L_1^2}\phi_h &= 0 \\
    \label{eq:partic}
    \ddxx \phi_p - \frac{1}{L_1^2}\phi_p &= -\frac{S}{D_1}
\end{align}
To solve for the particular solution, we could consider a wide range of functions whose differentials could somehow combine to produce the inhomogeneous term $-\frac{S}{D_1}$.
In this case, since the inhomogeneous term is a constant, we could first try simply a constant function $\phi_p(x)=A$:
\begin{align}
    \phi_p(x)&=A\\
    \ddxx A - \frac{1}{L_1^2}A &= -\frac{S}{D_1}\nonumber\\
    \frac{\Sigma_{a1}}{D_1}A &= \frac{S}{D_1}\nonumber\\
    A &= \frac{S}{\Sigma_{a1}}\\
    \label{eq:part-sol}
    \rightarrow\phi_p &= \frac{S}{\Sigma_{a1}}
\end{align}
Then, for the homogeneous expression in Eq.~\ref{eq:homo}, applying a Helmholtz equation substitution:
\begin{gather}
    \ddxx \phi_h - \frac{1}{L_1^2}\phi_h = 0 \nonumber\\
    \label{eq:homo-sol}
    \rightarrow\phi_h = C_1\sinh(\frac{x}{L_1}) + C_2\cosh(\frac{x}{L_1}) \\
\end{gather}
We can then combine the homogeneous (\ref{eq:homo-sol}) and particular (\ref{eq:part-sol}) solutions to find $\phi_1$:
\begin{align}
    \color{gray}\phi_1&\color{gray}=\phi_h+\phi_p\\
    \label{eq:reg1-source-sol}
    \phi_1 &= C_1\sinh(\frac{x}{L_1}) + C_2\cosh(\frac{x}{L_1}) + \frac{S}{\Sigma_{a1}} \\
    \label{eq:reg1-source-diff}
    \ddx \phi_1 &= \frac{1}{L_1}\left(C_1\cosh(\frac{x}{L_1}) + C_2\sinh(\frac{x}{L_1})\right)
\end{align}
Then, for region 2 (Eq.~\ref{eq:reg2}), which has no source and is thus homogeneous, a Helmholtz equation solution is much more quickly found:
\begin{gather}
    \ddxx\phi_2 - \frac{1}{L_2^2}\phi_2 = 0\nonumber\\
    \label{eq:reg2-sol}
    \phi_2 = C_3\sinh(\frac{x}{L_2}) + C_4\cosh(\frac{x}{L_2})\\
    \label{eq:reg2-diff}
    \ddx \phi_2 = \frac{1}{L_2}\left(C_3\cosh(\frac{x}{L_2}) + C_4\sinh(\frac{x}{L_2})\right)
\end{gather}
To solve for the four coefficients $C_1$ through $C_4$, we employ four boundary conditions: symmetry around no source at $x=0$, conservation of flux and current at $x=a$, and an extrapolated vacuum boundary at $x=a+b$ (ignoring the extrapolation distance $2D_2$).
Applying the boundary conditions at $x=0$ and $x=a+b$:
\begin{align}
    \label{eq:sym-bound}
    0 &= \ddx \phi_1(x=0) \\
      &= \frac{1}{L_1}\left(C_1\cosh(\frac{0}{L_1}) + C_2\sinh(\frac{0}{L_1})\right)\nonumber\\
      &= \frac{C_1}{L_1} = C_1\nonumber\\
    \label{eq:vac-bound}
    0 &= \phi_2(x=a+b)\\
      &= C_3\sinh(\frac{a+b}{L_2}) + C_4\cosh(\frac{a+b}{L_2})\nonumber\\
    \rightarrow\phi_2&=C_5\sinh(\frac{x-(a+b)}{L_2})
\end{align}
The remaining boundary conditions at $x=a$ yield a system of two equations for $C_2$ and $C_5$:
\begin{align}
    \label{eq:prob2-abounds}
    0 &= \phi_1(x=a) - \phi_2(x=a)\\
      &= C_2\cosh(\frac{a}{L_1}) + \frac{S}{\Sigma_{a1}} - C_3\sinh(\frac{a}{L_2}) - C_4\cosh(\frac{a}{L_2})\nonumber\\
    0 &= -D_1\ddx\phi_1(x=a) + D_2\ddx\phi_2(x=a) \\
      &= -\frac{D_1}{L_1}\left(C_2\sinh(\frac{a}{L_1})\right) + \frac{D_2}{L_2}\left(C_5\cosh(\frac{-b}{L_2})\right)\nonumber\\
\end{align}
Solving this system of equations yields:
\begin{align}
    \label{eq:prob2-coeffs}
    C_2 &= -\frac{D_2 L_1 S \cosh(\frac{b}{L_2})}{\left(\Sigma_{a1} \left(D_1 L_2 \sinh(\frac{a}{L_1}) \sinh(\frac{b}{L_2}) + D_2 L_1 \cosh(\frac{a}{L_1}) \cosh(\frac{b}{L_2})\right)\right)} \\
    C_5 &= -\frac{D_1 L_2 S \sinh(\frac{a}{L_1})}{\left(\Sigma_{a1} \left(D_1 L_2 \sinh(\frac{a}{L_1}) \sinh(\frac{b}{L_2}) + D_2 L_1 \cosh(\frac{a}{L_1}) \cosh(\frac{b}{L_2})\right)\right)}\nonumber
\end{align}
which correspond to flux solutions:
\begin{align}
    \label{eq:reduced-sols}
    \phi_1&=C_2\cosh(\frac{x}{L_1}) + \frac{S}{\Sigma_{a1}} \\
    \phi_2&=C_5\sinh(\frac{x-(a+b)}{L_2})\nonumber
\end{align}
This solution, for the given material and geometric properties, can be seen in \figref{P2_flux}, compared to varying numerical solutions found in the next section.
\subsection{Finite Difference Solution}
Consider again the diffusion equation as found in Eq.~\ref{eq:diffusion}:
\begin{equation*}
    \frac{1}{v}\frac{\partial \phi\rt}{\partial t}-\nabla\cdot D\nabla\phi\rt + \Sigma_a\rt\phi\rt = \nu_p\Sigma_f\rt\phi\rt + \sum_{i=1}^{N}\lambda_iC_i\rt + s\rt
\end{equation*}

To solve using a finite difference method, we replace each derivative with a finite difference approximation.
In an analytical solution where one would consider separate equations for each uniform region, then for each equation $D$ does not change with respect to position and $-\nabla D\nabla\phi\rt$ can be simplified to $-D\nabla^2\phi\rt$ (in the above 1D case, $-D\frac{\partial^2\phi\rt}{\partial x^2}$).
To use a finite difference method, with \emph{one} equation accounting for spatial variation in $D$ across the material, then this assumption must be withheld to find a \emph{conservative} finite difference approximation of $-\nabla D\nabla\phi\rt$:
% \newcommand{\lfrac}{\frac{D_{i-1}}{D_i+D_{i-1}}}
% \newcommand{\rfrac}{\frac{D_{i+1}}{D_i+D_{i+1}}}
% \begin{equation}
%     -\nabla D_i\nabla\phi_i \approx -\frac{2D_i}{h^2}\left(\right)
% \end{equation}
\newcommand{\pfrac}{\left(\frac{D_{i+1}}{D_i+D_{i+1}}\right)}
\newcommand{\mfrac}{\left(\frac{D_{i-1}}{D_i+D_{i-1}}\right)}
\begin{equation}
    \label{eq:conserv-stream}
    \left[-\ddx D\ddx\phi\right]_i = -\frac{D_i}{h^2/2}\left(\pfrac\left(\phi_{i+1}-\phi_i\right)+\mfrac\left(\phi_{i-1}-\phi_i\right)\right) 
\end{equation}
This approximation of $-\ddx D\ddx\phi$ can then be plugged into the non-multiplying diffusion equation to yield an equation true for a given node $i$:
\begin{equation}
    \label{eq:center-findiff-short}
    \left[-\ddx D\ddx\phi\right]_i + \Sigma_{a,i}\phi_i = S_i
\end{equation}

\newcommand{\dL}[1]{d_{L,#1}}
\newcommand{\dR}[1]{d_{R,#1}}
\newcommand{\dLR}[1]{d_{LR,#1}}
or, by expanding each term and substituting shorthand $\dL{i}=\frac{D_i}{h^2/2}\mfrac$ and $\dR{i}=\frac{D_i}{h^2/2}\pfrac$ for the repeated fractions in Eq.~\ref{eq:conserv-stream}:
\begin{equation}
    \label{eq:shorthand}
    \begin{split}
        \dL{i}&=\frac{D_i}{h^2/2}\mfrac\\
        \dR{i}&=\frac{D_i}{h^2/2}\pfrac
    \end{split}
\end{equation}
\begin{equation}
    \label{eq:center-findiff}
    -d_L\phi_{i-1} + \left(d_L+d_R+\Sigma_{a,i}\right)\phi_i - d_R\phi_{i+1} = S
\end{equation}
This equation applies to every of the $N+1$ nodes except for the leftmost $i=0$ and rightmost $i=N$, which lack such relative nodes $i-1$ and $i+1$, respectively.
Instead, at these points, we may apply the boundary conditions of $x=0$ and $x=a+b$.
At $x=0$, we apply the same symmetric boundary condition seen in Eq.~\ref{eq:sym-bound}, using a similar approach as in Eq.~\ref{eq:find_coeff} to make a right-handed approximation of $\ddx\phi$:
\begin{equation}
    \label{eq:sym-findiff}
    \ddx\phi_0 \approx \frac{3\phi_0-4\phi_1+\phi_2}{2h}=0
\end{equation}
Then, neglecting the extrapolated distance as in Eq.~\ref{eq:vac-bound}, no approximation of the derivative is necessary, and we simply set the final node's flux to zero:
\begin{equation}
    \label{eq:vac-findiff}
    \phi_N = 0
\end{equation}
We have now obtained a system of linear equations, where some linear combination of values of $\phi$ can be set to a known value.
These can be assembled into a matrix equation (where, for conciseness, $\dLR{i}=\dL{i}+\dR{i}$, and empty spaces are assumed to be 0):

\begin{align}
    \label{eq:prob2-matrix}
    \begin{bmatrix}
        -3 & 4 & -1 & & & \\
        -\dL{1} & \dLR{1}+\Sigma_{a,2} & -\dR{1} &&\\
        & -\dL{2} & \dLR{2}+\Sigma_{a,2} & -\dR{2} &\\
        &&\ddots& \ddots & \ddots\\
        &&-\dL{N-1} & \dLR{N-1}+\Sigma_{a,N-1} & -\dR{N-1}\\
        &&&&1
    \end{bmatrix}&
    \begin{bmatrix}
        \phi_0\\\phi_1\\\phi_2\\\vdots\\\phi_{N-1}\\\phi_N
    \end{bmatrix}
    &=&
    \begin{bmatrix}
        0\\S_1\\S_2\\\vdots\\S_{N-1}\\0
    \end{bmatrix} \\
    A\qquad&\qquad \phi&=&\qquad b\nonumber
\end{align}
This equation can then be used to solve for the flux distribution $\phi$ to be calculated simply by calculating matrices $A$ and $b$ and finding $\phi=A^{-1}b$.\\

Using Eq.~\ref{eq:prob2-matrix}, numerical solutions were obtained for simulations with cell counts $N\in\left\{5, 10, 20, 40\right\}$.
The flux distributions obtained for each simulation are displayed in \figref{P2_flux} and compared to the analytical solution for the same distribution.

\defplot{P2_flux}{The flux of the (1D, Cartesian, non-multiplying) as found by analytical and numerical solutions with varying numbers of cells $N$.}
\FloatBarrier

These solutions, as expected, seem to trend closer and closer to the analytical solution as $N$ increases.
With solutions using enough values of $N$, the change of error with respect to $h$ (in this case, equivalent to $\frac{a+b}{N}$) can be estimated.
If we expect the solution to have order $\bigO(h^n)$, then we can solve for $n$:
\begin{equation}
\label{eq:err-order}
\begin{split}
    \mathrm{error} &= Kh^n\\
    \log(\mathrm{error}) &= \log(K)+n\log(h)
\end{split}
\end{equation}
Taking a linear regression of $\log(\mathrm{error})$ and $\log(h)$ found using solutions with $N\in\{$25, 50, 100, 200, 400, 800, 1600$\}$ cells, an order of convergence was found of $\bigO(h^{1.00912})$, slightly better than first order.
The errors found for different values of $h$ may be seen on a log-log plot in \figref{P2_error}.\\

To estimate error, a relative L$^2$ norm\cite{novak_computer_2026} was computed for each flux solution $\phi$ and the exact analytical solution $\phi^*$:
\begin{equation}
    \label{eq:err-def}
    \mathrm{error} = \frac{\left|\left|\phi-\phi^*\right|\right|_2}{\left|\left|\phi^*\right|\right|_2}
\end{equation}
where the L$^2$ norm is defined for a vector $x$ of length $N+1$ in Eq.~\ref{eq:l2-def} below:
\begin{equation}
    \label{eq:l2-def}
    ||x||_2 = \sqrt{\sum_{i=1}^{N+1}x_i^2}
\end{equation}
\defplot{P2_error}{The error found in numerical simulations of cell counts $N\in\left\{25, 50, 100, 200, 400, 800, 1600\right\}$, compared to their cell sizes $h$. As expected, error increases with larger sizes $h$, doing so slightly faster than a first order error $\bigO(h^1)$.}
\FloatBarrier
\section{Problem 3: Multiplying Systems}

Consider again the neutron diffusion equation seen in Eq.~\ref{eq:diffusion}, this time simplified by assuming no external sources or delayed neutron production are present, but with considerable prompt fission:
\begin{equation}
    \label{eq:multiplying-diff}
    \frac{1}{v}\frac{\partial \phi}{\partial t} -\nabla\cdot D\nabla\phi\rt + \Sigma_a\rt\phi\rt = \nu\Sigma_f\rt\phi\rt
\end{equation}
In such multiplying systems, the time rate of change of $\phi$ is proportional to some spatial distribution of $\phi$ and nothing else, forming eigenvector combinations of flux distributions which exponentially increase or decay in time.
While a full and proper analysis might try to take into account every possible multiplying distribution $\phi$, each fundamental distribution, or ``mode'', corresponds to a different multiplication rate $k$, and thus for a system left to operate over long enough time, one can describe flux simply using the most dominant mode, that is, the one with largest multiplication rate $k$.\\

Thus suppose that for some flux distribution $\phi$, the distribution time change $\frac{1}{v}\frac{\partial \phi}{\partial t}$ is proportional to the current distribution of flux, such that on average every neutron released goes on to cause $k$ subsequent neutron releases through fission, a multiplication rate $k$.
Then, while the distribution itself would not be in steady state, that distribution, in a hypothetical material with $\nu=\frac{\nu}{k}$ would be in steady state, allowing us to write:
\begin{equation}
    \label{eq:pseudo-steady}
    -\nabla\cdot D\nabla\phi\rt + \Sigma_a\rt\phi\rt = \frac{1}{k}\nu\Sigma_f\rt\phi\rt
\end{equation}
and form some sort of Helmholtz equation to describe $\phi$.
Typically, this revolves around a Helmholtz coefficient termed ``material buckling'':
\begin{equation}
    \label{eq:mat-buckling}
    B=\frac{\nu\Sigma_f - \Sigma_a}{D}
\end{equation}
This ``material'' buckling can then be compared to the ``geometric'' buckling found through the geometry of the problem to determine what materials or what geometry would be required to produce a critical system (if such a feat is possible).\\

For a non-critical system of determined geometry and material, if the unknown multiplication rate $k$ is folded into this definition as in Eq.~\ref{eq:pseudo-steady}, one can define a hypothetically critical material buckling $B_c$:
\begin{equation}
    \label{eq:crit-buckling}
    B_c=\frac{\frac{1}{k}\nu\Sigma_f - \Sigma_a}{D}
\end{equation}

and form a Helmholtz equation for multiplying flux within a 1D medium of uniform properties:
\begin{equation}
    \label{eq:noncrit-mult}
    \begin{gathered}
        -D\nabla^2\phi + \Sigma_a\phi\rt = \frac{1}{k}\nu\Sigma_f\phi\rt
        \\
        -D\ddxx \phi(x) - \left(\frac{1}{k}\nu\Sigma_f-\Sigma_a\right)\phi(x) = 0 \\
        \ddxx\phi + \frac{\frac{1}{k}\nu\Sigma_f - \Sigma_a}{D}\phi = 0\\
        \ddxx\phi + B_c^2\phi = 0\\
        \phi = C_1\sin(B_cx) + C_2\cos(B_cx)
    \end{gathered}
\end{equation}
Using boundary conditions, an imposed scale (such as normalizing reactor power to 1), and assuming a non-trivial solution, one would then solve for a value of $B_c$ by solving for $B_c$ which allow the spatial distribution to line up with such conditions.
Because an individual value $B_c$ corresponds to an individual $k$ value (Eq.~\ref{eq:crit-buckling}), this condition typically reveals an infinite series of $B_c$.
However, since for systems after enough time we only need to consider the dominant mode of flux, we can consider only the smallest value of $B_c$, which results in the largest multiplication rate $k$:
\begin{equation}
    \label{eq:k-buck}
    k=\frac{\nu \Sigma_{f1}}{\left(\left(B_c^{2} D\right)+\Sigma_{a1}\right)}
\end{equation}
Then, one can describe the long-term flux distribution, using the actual material buckling $B$, and its multiplication rate $k$.
For a more thorough description of such combinations of modes of flux, one may wish to consult a separation-of-variables derivation of the time-dependent flux\cite{novak_lecture11_2026}.

\subsection{Analytical Solution}
Consider, for example, the system of \figref{prob2_setup}, but this time with a region 1 which has a fission cross-section $\Sigma_f$, rather than a distributed source.
In this problem, consider region~1 to have properties $\Sigma_{a1}=0.066$~[1/cm], $\Sigma_{f1}=0.02787$~[1/cm], $D_1=0.79$~[cm], $\nu=2.4$~[neutrons/fission], and region~2 the properties $\Sigma_{a2}=0.000709$~[1/cm] and $D_2=1.0$~[cm]\cite{novak_computer_2026}. \\

Region 1 may be described with the multiplying diffusion equation of Eq.~\ref{eq:noncrit-mult}:
\begin{equation}
    \begin{gathered}
        -D\nabla^2\phi_1 + \Sigma_a\phi_1\rt = \frac{1}{k}\nu\Sigma_f\phi_1\rt
        \\
        -D\ddxx \phi_1(x) - \left(\frac{1}{k}\nu\Sigma_f-\Sigma_a\right)\phi_1(x) = 0 \\
        \ddxx\phi_1 + \frac{\frac{1}{k}\nu\Sigma_f - \Sigma_a}{D}\phi_1 = 0\\
        \ddxx\phi_1 + B_c^2\phi_1 = 0\\
        \phi_1 = C_1\sin(B_cx) + C_2\cos(B_cx)\\
        \ddx\phi_1 = B_cC_1\cos(B_cx) + B_cC_2\sin(B_cx)
    \end{gathered}
\end{equation}
Using the symmetry boundary condition:
\begin{equation}
    \begin{aligned}
        \ddx\phi_1(0) = 0 &= B_cC_1\cos(0) + B_cC_2\sin(0)\\
        0&= B_cC_1\\
    \end{aligned}
\end{equation}
\begin{equation}
    \label{eq:reg1-mult-reduced}
    \begin{gathered}
        \color{gray}\rightarrow C_1=0\\
        \phi_1 = C_2\cos(B_cx)\\
        \ddx\phi_1 = -B_cC_2\sin(B_cx)
    \end{gathered}
\end{equation}
Region 2, which has neither distributed source nor multiplying material, can be described using an identical process as in Problem~2, Eqs.~\ref{eq:reg2-sol} to \ref{eq:reduced-sols}:
\begin{equation}
    \label{eq:reg2-reduced}
    \begin{aligned}
        \phi_2&=C_5\sinh(\frac{x-(a+b)}{L_2})\\
        \ddx\phi_2&= \frac{C_5}{L_2}\cosh(\frac{x-(a+b)}{L_2})
    \end{aligned}
\end{equation}

Combining the flux distributions in Eqs.~\ref{eq:reg1-mult-reduced} and \ref{eq:reg2-reduced} with the same boundary conditions at $x=a$ as fr the non-multiplying system (Eq.~\ref{eq:prob2-abounds}):
\begin{equation}
    \label{eq:prob3-abounds}
    \begin{aligned}
        0 &= \phi_1(x=a) - \phi_2(x=a)\\
        &= C_2\cos(B_ca) - C_5\sinh(\frac{-b}{L_2})\\
        0 &= -D_1\ddx\phi_1(x=a) + D_2\ddx\phi_2(x=a) \\
        &= D_1B_cC_2\sin(B_ca) + \frac{D_2C_5}{L_2}\cosh(\frac{-b}{L_2})\\
    \end{aligned}
\end{equation}
Since there is no constant term such as $S$ in this expression, trying to solve for this problem for $C_2$ and $C_5$ would only yield the trivial solution $C_2=C_5=0$, we instead use one boundary condition to solve for $C_5$ in terms of $C_2$:
\begin{equation}
    \label{eq:c5-c2}
    C_5 = C_2\frac{\cos(B_ca)}{\sinh(\frac{-b}{L_2})}
\end{equation}
and use the other as a condition on what $B_c$ must be for the system to be critical:
\begin{equation}
    \label{eq:crit-condition}
    \begin{aligned}
        D_1B_cC_2\sin(B_ca) &= -\frac{D_2C_5}{L_2}\cosh(\frac{-b}{L_2})\\
        D_1B_cC_2\sin(B_ca) &= -\frac{D_2C_2\frac{\cos(B_ca)}{\sinh(\frac{-b}{L_2})}}{L_2}\cosh(\frac{-b}{L_2})\\
        B_c\tan(B_ca) &= \frac{D_2}{D_1L_2\tanh(\frac{b}{L_2})}
    \end{aligned}
\end{equation}

Using a graphical tool like \hyperlink{https://www.desmos.com/calculator/uuepsi5ikz}{Desmos}, the values of $B_c$ for which this distribution is valid may be visualized, and the smallest $B_c$ found, as shown in \figref{P3_desmos}.

\defplot{P3_desmos}{For an example problem ($a=10$~cm, $b=50$~cm), the smallest value of $B_c$ which satisfies the criticality condition $B_c\tan(B_ca)= \frac{D_2}{D_1L_2\tanh(\frac{b}{L_2})}$}

Finally, to solve for $C_2$, some condition may be used to force a specific scale, such as by requiring the neutron production of the reactor to equal 1 (taking care to distinguish between the actual material buckling $B$ and the hypothetically critical buckling $B_c$):
\begin{equation}
    \label{eq:production-force}
    \begin{aligned}
        1=P&=\int_{0}^{a+b}\phi(x)\nu\Sigma_f(x)dx\\
        &=\int_{0}^{a}\phi_1(x)\nu\Sigma_{f,1}dx\\
        &=\int_{0}^{b}C_2\cos(Bx)\nu\Sigma_{f,1}dx\\
        &=\frac{C_2}{B}\sin(Bx)\nu\Sigma_{f,1}\\
        C_2&=\frac{B}{\sin(Bx)\nu\Sigma_{f,1}}
    \end{aligned}
\end{equation}
which may then be used alongside Eq.~\ref{eq:c5-c2} for the two final flux distributions:
\begin{equation}
    \label{eq:p3-dists}
    \begin{gathered}
        C_2=\frac{B}{\sin(Bx)\nu\Sigma_{f,1}}\\
        C_5=C_2\frac{\cos(B_ca)}{\sinh(\frac{-b}{L_2})}\\
        \phi_1= C_2\cos(B_cx)\\
        \phi_2=C_5\sinh(\frac{x-(a+b)}{L_2})
    \end{gathered}
\end{equation}

Returning to the criticality condition of Eq.~\ref{eq:crit-condition}, we may glean useful information about the reactor.\\

For instance, if we wanted a geometry which was perfectly critical, then we could impose that $B_c=B$, since $k=1$, and solve for $a$ or $b$ given the other:
\begin{equation}
    \label{eq:crit-ab}
    B\tan(Ba) = \frac{D_2}{D_1L_2\tanh(\frac{b}{L_2})}
\end{equation}
for instance, we could use this condition as reflector thickness $b$ approaches infinity to solve for the minimum thickness of the multiplying layer $a$:
\begin{equation}
    \label{eq:min-a}
    B\tan(Ba_{min}) = \frac{D_2}{D_1L_2}
\end{equation}
which, for the given material parameters, evaluates to 23.505~cm.
Likewise, the critical $a$ for a system with \emph{no} shielding would indicate the maximum thickness of the multiplying layer, above which no matter how little reflection was added, the reactor would be supercritical.
While Eq.~\ref{eq:crit-ab} is not itself defined at $b=0$, the same result may be analogously found by applying the (extrapolation-ignored extrapolated) vacuum boundary condition directly to $x=a$ in this scenario:
\begin{equation}
    \label{eq:max-a}
    \phi_1(a)=0=\cos\left(Ba\right) 
\end{equation}
which, for the given material parameters, evaluates to 46.852~cm.\\

Using Eq.~\ref{eq:crit-ab}, the reflector thickness $b$ required for several given lengths of $a$ to form critical reactors are calculated and tabulated in Table~\ref{tab:bvals} (where possible).

% Critical a when no multiplier: 46.85187 cm (\cos\left(Ba\right) = 0)
% Critical a as b approaches inf: 23.50502 cm (BD_{1} = \frac{D_{2}}{L_{2}\tan(Ba)})
\begin{table}
    \label{tab:bvals}
    \caption{The reflector thicknesses $b$ required to make a system with a multiplying region $a$ thick critical. Systems which with any thickness $b$ would be subcritical are denoted $k<1$, and ones which would be supercritical for any $b$ are denoted $k>1$.}
    \begin{tabular}{c|cccccccccc}
a (cm) & 5 & 10 & 15 & 20 & 25 & 30 & 35 & 40 & 45 & 50\\\hline
b (cm) & $k<1$ & $k<1$ & $k<1$ & $k<1$ & 56.187 & 28.302 & 16.901 & 8.997 & 2.350 & $k>1$
    \end{tabular}
\end{table}
\FloatBarrier
\subsection{Finite Difference Solution}
Similarly to as in the prior problem, the diffusion equation of Eq.~\ref{eq:pseudo-steady} may be formulated as a finite difference equation:
\begin{equation}
    \label{eq:mult-findiff}
    \left[-\ddx D\ddx\phi\right]_i + \Sigma_{a,i}\phi_i = \frac{1}{k}\nu\Sigma_{f,i}\phi_i
\end{equation}
with $\left[-\ddx D\ddx\phi\right]_i$ and shorthand terms $\dL{i}$ and $\dR{i}$ defined as in Eq.~\ref{eq:shorthand}:
\begin{equation}
    \label{eq:mult-findiff}
    -d_L\phi_{i-1} + \left(d_L+d_R+\Sigma_{a,i}\right)\phi_i - d_R\phi_{i+1} = \frac{1}{k}\nu\Sigma_{f,i}\phi_i
\end{equation}
For the boundary nodes $\phi_0$ and $\phi_N$, the same boundary conditions as in Eqs.~\ref{eq:sym-findiff} and \ref{eq:vac-findiff} may be applied.
Together, these yield a matrix equation similar to \ref{eq:prob2-matrix}

\begin{equation}
    \label{eq:prob3-matrix}
    \begin{aligned}
    \begin{bmatrix}
        -3 & 4 & -1 &  \\
        -\dL{1} & \dLR{1}+\Sigma_{a,2} & -\dR{1} &&\\
        &\ddots& \ddots & \ddots\\
        &-\dL{N-1} & \dLR{N-1}+\Sigma_{a,N-1} & -\dR{N-1}\\
        &&&1
    \end{bmatrix}&
    \begin{bmatrix}
        \phi_0\\\phi_1\\\vdots\\\phi_{N-1}\\\phi_N
    \end{bmatrix}
    &=&
    \frac{1}{k}
    \begin{bmatrix}
        0\\\nu_1\Sigma_{f,1}\\\vdots\\\nu_{N-1}\Sigma_{f,N-1}\\0
    \end{bmatrix} \\
    M\qquad&\qquad \phi&=&\qquad \frac{1}{k}F\phi
    \end{aligned}
\end{equation}
where $\frac{1}{k}F\phi$ expands to:
\begin{equation}
    \begin{split}
    \frac{1}{k}
    \begin{bmatrix}
        0 & &  & & & \\
        & \nu_1\Sigma_{f,1} & &&\\
        && \ddots &\\
        &&& \nu_{N-1}\Sigma_{f,N-1} &\\
        &&&&1
    \end{bmatrix}&
    \begin{bmatrix}
        \phi_0\\\phi_1\\\vdots\\\phi_{N-1}\\\phi_N
    \end{bmatrix}\\
    F\qquad&\qquad \phi
    \end{split}
\end{equation}
This equation is true for any (eigenvector) distribution of flux $\phi$ which is scaled every generation by an (eigenvalue) multiplication rate $k$, and the dominant (largest) value of $k$ and corresponding distribution $\phi$ may be solved for by manipulating the relationship to return a new $\phi_{n+1}$ using a previous guess of $k_n$ and $\phi_n$ (where $n$ represents the index of power iteration):
\begin{equation}
    \label{eq:pow-iter}
    \begin{gathered}
        M\phi = \frac{1}{k}F\phi \\
        \phi = M^{-1}\left(\frac{1}{k}F\phi\right)\\
        \phi_{n+1} = M^{-1}\left(\frac{1}{k_n}F\phi_n\right)\\
        k_{n+1} = \frac{\int_{0}^{a+b}\nu\Sigma_{f}\phi dx}{\int_{0}^{a+b}\Sigma_{a}\phi dx}
    \end{gathered}
\end{equation}
where the integral evaluations for $k_{n+1}$ are carried out using a numerical integral process such as a trapezoidal Reimann sum.
To maintain consistent scales of $\phi$, after every iteration the distribution is then normalized using a similar integral evaluation of power, to normalize power to 1:
\begin{gather}
    P=\int_{0}^{a+b}\Sigma_{f}\phi dx\\
    \phi = \frac{\phi}{P}
\end{gather}
\\

Starting from a reasonable initial guess of $k$ and $\phi$, this process may then be repeated until sufficient convergence criteria are met:
\begin{equation}
    \label{eq:convergence}
    \begin{gathered}
        \frac{||\phi_{n+1}-\phi_{n}||_2}{||\phi_n||_2} \leq 10^{-6}\\
        \frac{(k_{n+1}-k_n)^2}{(k_n)^2}\leq 10^{-6}
    \end{gathered}
\end{equation}
Using this method for discretizations of $N\in\left\{10, 20, 40, 80, 160\right\}$ cells, distributions of $\phi$ and their corresponding $k$ values were solved for, and may be seen in \figref{P3_a10} for geometry with $a=10$, in \figref{P3_a25} for $a=25$, and \figref{P3_a50} for $a=50$, with $k$ values additionally compared in Table~\ref{tab:kvals}. 
These distributions and $k$ values are also compared to those found using the previously described analytical solution.
\FloatBarrier

% K results 
% a (cm) & b (cm) & Numerical Solutions &&&&& Analytical
% 10 & 50 & 0.96688 & 0.97559 & 0.97257 & 0.97407 & 0.97331 & 0.97356\\
% 25 & 50 & 1.00042 & 0.99925 & 0.99981 & 0.99951 & 0.99966 & 0.9960\\
% 50 & 50 & 1.00858 & 1.00828 & 1.00812 & 1.00804 & 1.00800 & 1.00796\\
\defplot{P3_a10}{The normalized flux distributions of geometry with $a=10$ for numerical solutions with $N\in\left\{10, 20, 40, 80, 160\right\}$ cells, compared to an analytical solution.}
\defplot{P3_a25}{The normalized flux distributions of geometry with $a=25$ for numerical solutions with $N\in\left\{10, 20, 40, 80, 160\right\}$ cells, compared to an analytical solution.}
\defplot{P3_a50}{The normalized flux distributions of geometry with $a=50$ for numerical solutions with $N\in\left\{10, 20, 40, 80, 160\right\}$ cells, compared to an analytical solution.}

\begin{table}
    \label{tab:kvals}
    \caption{The k values found using finite difference numerical simulations with different numbers $N$ of cells, compared to those found analytically, for geometries of different lengths $a$ and $b$.}
    \begin{tabular}{cc|cccccc}
&Geometry& k \\
a (cm)&b (cm)& N = 10  & N = 20  & N = 40  & N = 80  & N = 160 & Analytic\\\hline
10 & 50 & 0.96688 & 0.97559 & 0.97257 & 0.97407 & 0.97331 & 0.97356 \\
25 & 50 & 1.00042 & 0.99925 & 0.99981 & 0.99951 & 0.99966 & 0.99961 \\
50 & 50 & 1.00858 & 1.00828 & 1.00812 & 1.00804 & 1.00800 & 1.00796 \\
    \end{tabular}
\end{table}
\FloatBarrier
\section{Conclusions}\label{sec:conc}

The objectives of this report were largely met.
Finite-difference simulations of both the multiplying and non-multiplying systems were successfully constructed and observed to converge towards the flux distributions and (for the multiplying system) multiplication rates predicted using an exact analytical solution.\\

Namely, the convergence of order $\bigO(h^{1.009})$ observed for the numerical simulation of the non-multiplying system is sufficiently high as expected.
In addition, several properties about the general design of the multiplying system were assessed using the analytical solution, such as minimum and maximum critical region~1 thicknesses $a$ (23.505 and 46.852~cm, respectively) and the reflector thickness $b$ required for arbitrary multiplying region thicknesses $a$.\\

The effective convergence of all numerical solutions to exact analytical equivalents highlights both their robustness where analytical solutions may be problematic or painful, and their weakness where computational resources are limited.
Additionally, the finite-difference solutions used are still designed around simple, 1D nodal geometry, and would be ill-suited towards finely detailed geometry themselves.\\

These results give much opportunity to expand on these simulations in future work, such as by discretizing over more complicated or multidimensional geometry, solving for flux over time, or discretizing over energy (multigroup theory).

% \bibliographystyle{sty/ieeetran}
\pagebreak\printbibliography

\appendix
\pagebreak

\end{document}
